\section{Taller resuelto: Trim no lineal con empuje en cuerpo}

\noindent\textbf{Nivel:} Intermedio.

\apartado{1. Enunciado}{Resuelva el equilibrio de un avión a 66 m/s y rho=1.04 con empuje paralelo a cuerpo. Masa 2450 kg, S=26 m², CL=0.24+5.2alpha+0.28delta, Cm=0.045-0.72alpha-1.08delta y polar 0.026+0.046CL².}

\apartado{2. Datos}{Vuelo horizontal: cabeceo theta=alpha; ángulos en radianes; límites de elevador ±0.35 rad y T entre 0 y 5000 N.}

\apartado{3. Incógnitas}{Las magnitudes y comprobaciones solicitadas en el caso.}

\apartado{4. Sistema de referencia}{Ejes viento para fuerzas y cuerpo para dirección de empuje.}

\apartado{5. Hipótesis físicas}{Modelo lineal de coeficientes, geometría exacta del empuje y momento propulsivo nulo respecto al CG.}

\apartado{6. Modelo matemático}{Trim no lineal con empuje en cuerpo}

\apartado{7. Ecuaciones simbólicas}{Se parte de los balances y definiciones siguientes.}

\begin{equation}T\cos\alpha-D=0\end{equation}

\begin{equation}L+T\sin\alpha-W=0\end{equation}

\begin{equation}C_{m0}+C_{m\alpha}\alpha+C_{m\delta}\delta=0\end{equation}

\apartado{8. Desarrollo algebraico}{Las operaciones se realizan en este orden, conservando las condiciones del modelo.}

\begin{enumerate}

\item La fuerza de empuje se proyecta en dos direcciones.

\item Se plantean tres incógnitas y tres residuos; la incidencia aparece además dentro de CL.

\item Se resuelve con una estimación inicial física y se comprueban límites de mando y empuje.

\item Se sustituyen resultados sin redondear en los tres balances.

\end{enumerate}

\apartado{9. Sustitución numérica}{Se conservan los resultados intermedios sin redondear y se muestran todas las operaciones escalares disponibles.}

\textbf{Resolución numérica que puede reproducirse.} La ecuación de momento se despeja como $\delta=(0.045-0.72\alpha)/1.08$. El equilibrio tangencial da $T=D/\cos\alpha$. Sustituir ambos en el vertical reduce tres incógnitas a una: $f(\alpha)=qSC_L(\alpha)+D(\alpha)\tan\alpha-mg_0=0$, con $qS=58893.12\,\mathrm N$, $C_L=0.24+5.2\alpha+0.28\delta$ y $D=qS(0.026+0.046C_L^2)$. Se elige un intervalo con cambio de signo. En cada paso $c=(a+b)/2$ y se conserva la mitad que mantiene signos opuestos. El ancho final acota el error en la raíz, bajo continuidad.
\begin{center}\begin{tabular}{rrrrr}
\toprule Paso & $a$ & $b$ & $c$ & $f(c)$\\\midrule
$1$ & $0$ & $0.2$ & $0.1$ & $20627.98$\\
$2$ & $0$ & $0.1$ & $0.05$ & $5668.519$\\
$3$ & $0$ & $0.05$ & $0.025$ & $-1775.67$\\
$4$ & $0.025$ & $0.05$ & $0.0375$ & $1944.144$\\
$8$ & $0.0296875$ & $0.03125$ & $0.03046875$ & $-148.7602$\\
$16$ & $0.03096313$ & $0.03096924$ & $0.03096619$ & $-0.7388058$\\
$24$ & $0.03096867$ & $0.03096869$ & $0.03096868$ & $0.002594401$\\
$35$ & $0.03096867$ & $0.03096867$ & $0.03096867$ & $1.425818\times10^{-6}$\\
\bottomrule\end{tabular}\end{center}
El resultado queda en $[0.03096866928,0.03096866928]\,\mathrm{rad}$. Después se recuperan $\delta$ y $T$ con las dos relaciones despejadas, se comprueban sus límites y se sustituyen en los tres residuos. No se impone $L=W$ porque el empuje tiene componente vertical.





\apartado{10. Resultados con unidades}{Los resultados pertenecen al modelo y al intervalo de datos declarados.}

\begin{center}\begin{tabular}{lr}\toprule Magnitud&Resultado\\\midrule

Incidencia & $0.03096867\,\mathrm{rad}$\\

Elevador & $0.02102089\,\mathrm{rad}$\\

Empuje & $1980.758\,\mathrm N$\\

Residuo tangencial & $4.547474\times10^{-13}\,\mathrm N$\\

Residuo normal & $3.637979\times10^{-12}\,\mathrm N$\\

Residuo de momento adimensional & $3.469447\times10^{-18}\,$\\

\bottomrule\end{tabular}\end{center}

\apartado{11. Interpretación}{El componente vertical del empuje reduce ligeramente la sustentación aerodinámica necesaria frente a L=W.}

\apartado{12. Comprobación independiente y alcance}{Resolver desde varias estimaciones iniciales y comprobar residuos por componente; compararlo con el modelo que omite T sin(alpha).}

\apartado{13. Limitaciones}{La norma combina un residuo adimensional de momento con fuerzas sólo como diagnóstico de raíz; una tolerancia física debe escalar ambos tipos por separado.}

\par\smallskip\noindent{\footnotesize\textit{Datos y solución originales. Las sustituciones registradas se reproducen con SymPy y Studium; eso no establece validación experimental ni revisión experta.}}
